% Глава: 12-processes
% Статус: наполнено (первый прогон).
\chapter{Процессы и ресурсы}
\label{ch:12-processes}

\section{Что такое процесс}

\textbf{Процесс} --- это программа в состоянии исполнения:
свой набор инструкций, память, открытые файлы, потоки.
Каждый процесс получает уникальный идентификатор \textbf{PID}
(Process ID) и работает от имени пользователя
(\nameref{ch:05-permissions}).

В этой главе процессы --- это то, что запущено на вашей ВМ:
компиляторы, тесты, контейнеры на сервере \fpath{pglab.nsu.ru}.

Процесс имеет \textbf{состояния} (видны в \cmd{ps}):
\begin{itemize}
  \item \cmd{R} --- running / runnable (выполняется или ждёт CPU);
  \item \cmd{S} --- sleeping (ждёт событие: ввод, диск, сеть);
  \item \cmd{Z} --- zombie (завершился, но родитель ещё не
    «забрал» результат).
\end{itemize}

\section{Иерархия процессов}

Все процессы, кроме первого (\cmd{init} / \cmd{systemd}),
являются чьими-то дочерними. Родительский идентификатор ---
\textbf{PPID}. Когда оболочка запускает команду
(\nameref{ch:03-shell}), команда --- дочерний процесс
оболочки.

\begin{figure}[htb]
  \centering
  \begin{tikzpicture}[node distance=1.3cm]
    \node[process] (init) {systemd (PID 1)};
    \node[process, below left=1.2cm and -0.4cm of init] (sshd) {sshd};
    \node[process, below right=1.2cm and -0.4cm of init] (cron) {cron};
    \node[process, below left=1.2cm and -0.4cm of sshd] (bash) {bash};
    \node[process, below right=1.2cm and -0.4cm of sshd] (vim) {vim};
    \draw[flow] (init) -- (sshd);
    \draw[flow] (init) -- (cron);
    \draw[flow] (sshd) -- (bash);
    \draw[flow] (sshd) -- (vim);
  \end{tikzpicture}
  \caption{Дерево процессов: каждый процесс имеет родителя.}
  \label{fig:proc-tree}
\end{figure}

\section{Просмотр процессов}

\begin{minted}[bgcolor=gray!6]{bash}
ps aux                 # все процессы: пользовательские и системные
ps -ef                 # полный формат, с PPID
ps -u <user>           # процессы одного пользователя
ps aux | grep <name>   # найти процесс по именем
top                    # «живой» просмотр с обновлением
htop                   # интерактивно, с цветом, удобнее читать
\end{minted}

Колонки \cmd{ps aux}, которые стоит читать: \textbf{USER},
\textbf{PID}, \textbf{\%CPU}, \textbf{\%MEM},
\textbf{COMMAND}. В \cmd{top}/\cmd{htop} сверху видна
сводка: загрузка CPU, свободная память, количество
процессов.

\section{Сигналы и завершение процессов}

Процессу можно отправить \textbf{сигнал} --- системное
уведомление. Важнейшие:

\begin{center}
\begin{tabular}{clp{0.5\linewidth}}
\toprule
\textbf{Сигнал} & \textbf{Номер} & \textbf{Действие} \\
\midrule
SIGHUP  & 1  & «отключён терминал»; процесс обычно завершается \\
SIGINT  & 2  & прерывание (обычно \cmd{Ctrl+C} в терминале) \\
SIGQUIT & 3  & завершение с дампом (редко) \\
SIGTERM & 15 & вежливая просьба завершиться (по умолчанию) \\
SIGKILL & 9  & принудительное завершение (процесс не может перехватить) \\
\bottomrule
\end{tabular}
\end{center}

\begin{minted}[bgcolor=gray!6]{bash}
kill 1234              # SIGTERM для PID 1234
kill -9 1234           # SIGKILL, принудительно
kill -TERM <pid>       # явно
pkill <name>           # по имени процесса
pkill -f "python app"  # по полной командной строке
\end{minted}

Правило: сначала \cmd{SIGTERM}, и только если процесс не
реагирует --- \cmd{SIGKILL}. Вежливое завершение даёт
процессу закрыть файлы и освободить ресурсы.

\section{Фоновые и передние задачи}

Запустить процесс в фоне --- добавить \cmd{\&} в конец
команды:
\begin{minted}[bgcolor=gray!6]{bash}
long_job.sh &        # запустить в фоне
jobs               # фоновые задания этой оболочки
fg                 # вернуть последнее задание на передний план
fg %2              # вернуть задание №2 на передний план
bg                 # возобновить остановленное задание в фоне
Ctrl+Z             # остановить текущее задание
\end{minted}

\section{Ресурсы: CPU и память}

Linux выделяет процессу ресурсы по запросу: долю CPU
(время процессора) и виртуальную память
(\nameref{ch:16-virtual-memory}). Мониторинг:
\begin{minted}[bgcolor=gray!6]{bash}
top                  # загрузка по CPU и по процессам
htop                 # интерактивно
vmstat 1             # системная статистика раз в секунду
free -h              # сводка по памяти
\end{minted}

\section{Типичные ошибки}
\begin{itemize}
  \item Процесс «завис» --- посмотрите в \cmd{top}, не
    использует ли он 100\% CPU; иногда достаточно
    \cmd{kill} (SIGTERM).
  \item \cmd{kill 9} не работает --- правильно
    \cmd{kill -9}; без знака минус это PID 9, а не флаг.
  \item Задача в фоне «потерялась» --- \cmd{jobs},
    затем \cmd{fg} или \cmd{bg}.
  \item После \cmd{Ctrl+C} сервер не запустить снова ---
    убедитесь, что старый процесс реально завершён
    (\cmd{ps aux | grep ...}).
\end{itemize}

\section{Связь с остальным курсом}

Процессы --- центральное понятие ОС. Потоки вывода
(\nameref{ch:07-text-pipelines}) принадлежат процессам;
запуск и остановка контейнеров
(\nameref{ch:19-docker-basics}) --- это запуск и
остановка процессов в изолированном окружении.
Взаимодействие процессов между собой ---
\nameref{ch:18-ipc}; распределение памяти ---
\nameref{ch:16-virtual-memory}.

\section{Задачи для самопроверки}
\begin{enumerate}
  \item Как узнать PID запущенного процесса?
  \item Чем \cmd{SIGTERM} отличается от \cmd{SIGKILL}?
    Почему второй «невозможно перехватить»?
  \item Запустите долгую команду в фоне, остановите её
    (\cmd{Ctrl+Z}), затем возобновите в фоне.
  \item Как увидеть, какие процессы принадлежат текущему
    пользователю?
  \item Что такое zombie-процесс и почему он возникает?
\end{enumerate}
